\FloatBarrier
\section{High-Level Language Design}
\label{sec:language}
\FloatBarrier
\subsection{Dialects, Operations, and Attributes}
\label{sec:dialects}
A dialect groups user-defined operations with their operand signatures,
attribute schemas, and semantics. Users define an operation once and import
it into multiple rules; different dialects can provide different operation
vocabularies. Attributes configure an operation independently of its tensor
operands. For example, a transpose carries a permutation list. The following
excerpt declares its signature and attribute in a dialect file,
\code{tensor.tepl}:
\begin{codeblock}[H]
\begin{lstlisting}[language=TEPL]
dialect TensorLang {
  op transpose(input: tensor) -> tensor {
    attrs {
      permutation: index[];
    }
  }
}
\end{lstlisting}
\end{codeblock}
This excerpt introduces the schema; it omits the semantic blocks. In a
semantic definition, \code{dtype} specifies element-type behavior,
\code{axes} declares coordinate bindings, legality, and output order, and
\code{semantics} gives the value at those coordinates. For transpose,
\code{dtype} preserves the element type. Transpose changes output coordinate
order while preserving the coordinate-indexed input value, so a value clause
\code{yield input} does not imply an unchanged tensor interface. We defer
the representation and composition of these blocks to Axes IR in
Section~\ref{sec:axes-ir}.

\FloatBarrier
\subsection{Imports and Rewrite Rules}
A rule imports a dialect and selects the operations it uses:
\begin{codeblock}[H]
\begin{lstlisting}[language=TEPL]
from "tensor.tepl" import TensorLang as t;
use t::{transpose};
\end{lstlisting}
\end{codeblock}
Rewrites are written as algebraic equations with an LHS pattern followed by an
RHS replacement. Nested operations denote producer--consumer edges, while
tensor captures identify matched values or subcomputations. Repeated captures
refer to the same matched value. This application graph is distinct from the
semantic graph used to reason about tensor elements and their coordinates.

Our running pattern is
\code{(transpose[@outer] (transpose[@inner] X))}.
Its intended replacement is a single transpose of \code{X}; the remaining
task is to construct that transpose's permutation. Tensor declarations may
also constrain shape and dtype, but equal dimension sizes alone do not
establish coordinate identity.

\FloatBarrier
\subsection{Attribute Binding and Reuse}
\label{sec:binding}
The annotation \code{op[@name]} binds an operation's attribute descriptor on
the LHS. If LHS attributes are omitted, the matcher captures each operation's
descriptor independently and anonymously. A name is needed when an expression
or an RHS operation refers to that descriptor. Captures are available in
\code{where} conditions, \code{derive} expressions, and RHS operations.
Field access such as \code{@outer.permutation} selects an individual
attribute where the expression context supports it.

\begin{figuretext}{figures/Attribute Binding and Derivation.pdf}{Captured operation descriptors can be read and transformed to configure replacement operations.}{fig:attribute-binding}{An input graph with transpose and dot is replaced by a dot. Colored descriptor
bindings connect the rule to a table with two captured rows and one derived row.}
Reusing \code{@name} on a compatible RHS operation reuses its descriptor;
a derived descriptor can configure an RHS operation in the same way.
Compatibility checks the attribute structure, rather than the mathematical
correctness of that choice. In the transpose-composition example, binding
\code{@inner} and \code{@outer} makes both permutations available, but copying
either one alone generally does not give the replacement permutation.

Figure~\ref{fig:attribute-binding} illustrates capture and derivation through
a separate transpose-absorption example. The input graph supplies \code{@tr}
and \code{@d}; a user-written \code{f_transform} computes \code{@new_d}.
The table distinguishes the two captured descriptors from the derived one.
Its short fields \code{perm}, \code{l_dim}, and \code{r_dim} belong to a
schematic dialect, rather than the generalized-dot schema used later.
For the displayed permutation \code{[1,0,2]}, the left contraction position
changes from \code{1} in the transposed operand to \code{0} in the original
operand. The figure illustrates descriptor dataflow for this instance; an
arbitrary transpose-absorption rule still requires semantic checking.
\end{figuretext}

\FloatBarrier
\subsection{Writing Attribute Changes}
Conditions and replacement attributes are metadata computations. Expressions
can read captured attributes and available shape and dtype bindings, but they
do not execute tensor computations. Local bindings, builtins, lists, indexing,
conditionals, and comprehensions provide the computations needed to build
descriptors. For transpose composition,
\code{gather(@inner.permutation, @outer.permutation)} computes the replacement
permutation by selecting entries of the inner permutation in outer-permutation
order. Descriptor constructors collect computed fields into immutable
attributes for an RHS operation. Runtime expressibility and semantic
checkability are distinct: some expressions can execute even when the
compiler cannot establish that they preserve a rule's tensor result.

\FloatBarrier
\subsection{Conditions with \texorpdfstring{\code{where}}{where}}
\label{sec:where}
A \code{where} block ends with the condition deciding whether a rule applies.
Any local calculations precede that condition. Only \code{let} bindings end
with a semicolon. Conditions can use captured metadata, builtins, and local
expressions, combining predicates with \code{&&} and \code{||}. For the running
transpose pattern, an illustrative clause is:
\begin{codeblock}[H]
\begin{lstlisting}[language=TEPL]
where {
  let rank = len(@inner.permutation);
  len(@outer.permutation) == rank
}
\end{lstlisting}
\end{codeblock}
The length test illustrates field access and local computation. For a legal
composition of transposes, the equality already follows from operation
legality. A rule can omit \code{where} when no additional condition is needed.

Conditions may express semantic applicability or application policy; TEPL does
not infer profitability decisions. They can also state axis-routing
relationships, for example:
\begin{codeblock}[H]
\begin{lstlisting}[language=TEPL]
@outer.lhs_batching_dimensions == @inner->B
\end{lstlisting}
\end{codeblock}
Its mapping interpretation is introduced in Section~\ref{sec:partitioning}. For a verified semantic rule,
every correctness obligation must follow from operation legality and supported
user-written conditions.

\FloatBarrier
\subsection{Attribute Construction with \texorpdfstring{\code{derive}}{derive}}
The \code{derive} block names replacement descriptors computed from captured
information. Transpose composition combines two captures and an explicit
descriptor constructor:
\begin{codeblock}[H]
\begin{lstlisting}[language=TEPL]
rule compose_transposes_explicit {
  (transpose[@outer] (transpose[@inner] X))
  => (transpose[@combined] X)
  derive {
    @combined = {
      permutation: gather(@inner.permutation, @outer.permutation)
    };
  }
}
\end{lstlisting}
\end{codeblock}
If the inner permutation is $p$ and the outer permutation is $q$, the resulting
permutation is $r[i]=p[q[i]]$. The three parts of a rule have distinct roles:
\code{where} decides applicability, \code{derive} constructs attributes, and
the RHS specifies the replacement computation.

Local bindings can name intermediate metadata computations. A spread such as
\code{{ ..@outer, permutation: permutation }} copies fields before applying
explicit overrides. The \code{where} and \code{derive} blocks have separate
local scopes; derivations can read captured descriptors and earlier derived
descriptors. With the transpose semantics, this explicit constructor is within
the supported semantic-verification fragment. Checking its attribute types
alone would not establish that the rewrite preserves the result.

\FloatBarrier
\subsection{Automatic Attributes with \texorpdfstring{\code{auto}}{auto}}
\label{sec:auto}
Users can omit replacement attributes when operation semantics support
synthesis. An explicit \code{[auto]} annotation requests the same inference:
\begin{codeblock}[H]
\begin{lstlisting}[language=TEPL]
rule compose_transposes_auto {
  (transpose (transpose X)) => (transpose X)
}
\end{lstlisting}
\end{codeblock}
\begin{figuretext}{figures/Introducing auto.pdf}{\code{auto} becomes a reusable attribute function at compilation and is evaluated on matched metadata at runtime.}{fig:introducing-auto}{A vertical flow separates the user rule and operation semantics, compiler
synthesis and verification, and runtime application of the generated function.}
Matching supplies tensor captures and LHS descriptors. Semantic compilation
then derives a reusable program for the RHS permutation. The explicit rule
specifies how to combine two permutations; the automatic rule asks the compiler
to derive that computation from transpose semantics, again obtaining
$r[i]=p[q[i]]$.

Each omitted RHS descriptor is inferred independently. A named annotation
\code{[@name]} instead reuses an existing descriptor, while
\code{[auto, field = value]} pins one field during inference. Users still
provide operation semantics and any required correctness conditions.
Unsupported inference is rejected with a diagnostic. Constructor synthesis
occurs when the rule is compiled; application evaluates the completed
constructor on the matched metadata. Completed constructors can execute in
the greedy rewrite engine or in generated Rust with the shared runtime.
C++ generation remains unavailable.

Figure~\ref{fig:introducing-auto} makes the compile-time/runtime boundary
explicit using matrix transpose absorption. The compiler synthesizes and
checks a function $f(p,q)$ from the rule and operation semantics. Each runtime
match supplies the transpose permutation $p$ and dot descriptor $q$, and the
function supplies the replacement dot descriptor. In the matrix instance,
absorbing permutation \code{[1,0]} changes the left contraction position from
$0$ to $1$. The figure shows only that changing field; it does not imply that
other descriptor fields can be ignored in a complete operation definition.
\end{figuretext}

\FloatBarrier
\subsection{From \texorpdfstring{\code{auto}}{auto} to Attribute Synthesis}
The \code{auto} annotation lets users leave RHS attributes to the compiler,
but the compiler must still determine how to construct them. Matching exposes
known attributes and tensor metadata from the LHS. The challenge is to turn
that information into a program that computes the missing RHS attributes
across ranks and shapes and preserves the computation under the user's
conditions. The next section defines this problem and its correctness contract.
